%This chapter was modified on 4/2/97.
%\setcounter{chapter}{6}
\chapter[Jumlah Peubah Acak]{Jumlah Peubah Acak yang Saling Bebas}\label{chp 7}  

\section{Jumlah Peubah Acak Diskret}\label{sec 7.1} 


Dalam bab ini kita beralih ke pertanyaan penting mengenai cara menentukan distribusi
jumlah peubah acak yang saling bebas berdasarkan distribusi masing-masing peubah
penyusunnya.  Dalam bagian ini kita hanya meninjau jumlah peubah acak diskret; kasus
peubah acak kontinu akan dibahas pada bagian berikutnya.

Di sini kita hanya meninjau peubah acak yang nilainya berupa bilangan bulat.  Dengan
demikian, fungsi distribusinya didefinisikan pada bilangan-bilangan bulat tersebut.  Akan
lebih mudah jika kita menganggap bahwa fungsi distribusi ini didefinisikan untuk
\emx {semua} bilangan bulat, dengan menetapkan nilainya sama dengan~0 jika sebelumnya
tidak didefinisikan.

\subsection*{Konvolusi}

\par
Misalkan $X$ dan $Y$ dua peubah acak diskret yang saling bebas dengan fungsi distribusi
$m_1(x)$ dan $m_2(x)$.  Misalkan $Z = X+Y$.  Kita ingin menentukan fungsi distribusi
$m_3(x)$ dari $Z$.  Untuk melakukannya, cukup tentukan peluang bahwa $Z$ bernilai $z$,
dengan $z$ sebarang bilangan bulat.  Misalkan $X = k$, dengan $k$ suatu bilangan bulat.
Maka $Z = z$ jika dan hanya jika $Y = z-k$.  Jadi, kejadian $Z = z$ merupakan gabungan
kejadian-kejadian yang saling lepas berpasangan
$$(X = k)\  \mbox{dan\ }(Y = z-k)\ ,$$
dengan $k$ menjelajahi seluruh bilangan bulat.  Karena kejadian-kejadian ini saling lepas
berpasangan, kita memiliki
$$P(Z = z) = \sum_{k = -\infty}^\infty P(X = k)\cdot P(Y = z - k)\ .$$
Dengan demikian, kita telah memperoleh fungsi distribusi peubah acak $Z$.  Hal ini
mengantarkan kita pada definisi berikut.

\begin{definition}\label{defn 7.1}
Misalkan $X$ dan $Y$ dua peubah acak bernilai bilangan bulat yang saling bebas, dengan
fungsi distribusi masing-masing $m_1(x)$~dan~$m_2(x)$.  Maka {\em
konvolusi}\index{konvolusi} dari
$m_1(x)$~dan~$m_2(x)$ adalah fungsi distribusi $m_3 = m_1*m_2$ yang diberikan oleh
$$ m_3(j) = \sum_k m_1(k) \cdot m_2(j - k)\ ,$$
untuk $j = \ldots,\ -2,\ -1,\ 0,\ 1,\ 2,\ \ldots$.  Fungsi $m_3(x)$ merupakan fungsi
distribusi peubah acak $Z = X + Y$.
\end{definition}

\par
Mudah dilihat bahwa operasi konvolusi bersifat komutatif, dan mudah pula ditunjukkan
bahwa operasi tersebut juga bersifat asosiatif.
\par
Sekarang misalkan $S_n = X_1 + X_2 +\cdots+ X_n$ merupakan jumlah $n$ peubah acak yang
saling bebas dari suatu proses percobaan saling bebas, dengan fungsi distribusi yang sama,
yaitu $m$, yang didefinisikan pada bilangan bulat.  Maka fungsi distribusi $S_1$ adalah $m$.
Kita dapat menulis
$$
S_n = S_{n - 1} + X_n\ .
$$
Jadi, karena kita mengetahui bahwa fungsi distribusi $X_n$ adalah $m$, kita dapat
menentukan fungsi distribusi $S_n$ dengan induksi.

\begin{example}
Sebuah dadu dilempar dua kali.  Misalkan $X_1$ dan $X_2$ adalah hasilnya, dan misalkan
$S_2 = X_1 + X_2$ adalah jumlah kedua hasil tersebut.  Maka $X_1$ dan $X_2$ memiliki
fungsi distribusi yang sama:
$$
m = \pmatrix{
1 & 2 & 3 & 4 & 5 & 6 \cr
1/6 & 1/6 & 1/6 & 1/6 & 1/6 & 1/6\cr}.
$$
Fungsi distribusi $S_2$ kemudian merupakan konvolusi distribusi ini dengan dirinya
sendiri.  Jadi,

\begin{eqnarray*}
P(S_2 = 2) &=& m(1)m(1) \\
           &=& \frac 16 \cdot \frac 16 = \frac 1{36}\ , \\
P(S_2 = 3) &=& m(1)m(2) + m(2)m(1) \\
           &=& \frac 16 \cdot \frac 16 + \frac 16 \cdot \frac 16 = \frac 2{36}\ ,\\
P(S_2 = 4) &=& m(1)m(3) + m(2)m(2) + m(3)m(1) \\
           &=& \frac 16 \cdot \frac 16 + \frac 16 \cdot \frac 16 + \frac 16
\cdot \frac 16 = \frac 3{36}\ .\\
\end{eqnarray*}
Dengan melanjutkan cara ini, kita akan memperoleh $P(S_2 = 5) = 4/36$, $P(S_2 = 6) = 5/36$,
$P(S_2 = 7) = 6/36$, $P(S_2 = 8) = 5/36$, $P(S_2 = 9) = 4/36$, $P(S_2 = 10) =
3/36$, $P(S_2 = 11) = 2/36$, dan $P(S_2 = 12) = 1/36$.

Distribusi $S_3$ kemudian merupakan konvolusi distribusi $S_2$ dengan distribusi $X_3$.
Jadi,
\begin{eqnarray*}
P(S_3 = 3) &=& P(S_2 = 2)P(X_3 = 1) \\
           &=& \frac 1{36} \cdot \frac 16 = \frac 1{216}\ , \\
P(S_3 = 4) &=& P(S_2 = 3)P(X_3 = 1) + P(S_2 = 2)P(X_3 = 2) \\
           &=& \frac 2{36} \cdot \frac 16 + \frac 1{36} \cdot \frac 16 = \frac
3{216}\ ,\\
\end{eqnarray*}
dan seterusnya.
\par
Pekerjaan ini jelas membosankan, sehingga sebaiknya ditulis sebuah program untuk melakukan
perhitungan tersebut.  Untuk itu, mula-mula kita menulis program yang membentuk konvolusi
dua densitas $p$~dan~$q$ serta mengembalikan densitas $r$.  Selanjutnya, kita dapat menulis
program untuk mencari densitas jumlah $S_n$ dari $n$ peubah acak yang saling bebas dengan
densitas yang sama, yaitu $p$, setidaknya dalam kasus ketika peubah acak tersebut memiliki
sejumlah berhingga nilai yang mungkin.
\putfig{3.5truein}{PSfig7-1}{Densitas $S_n$ untuk pelemparan dadu $n$ kali.}{fig 7.1} 
\par
Menjalankan program ini untuk contoh pelemparan dadu sebanyak $n$ kali dengan~$n = 10,\ 20,\ 30$
menghasilkan distribusi yang ditunjukkan pada Gambar~\ref{fig 7.1}.  Terlihat
bahwa, seperti dalam percobaan Bernoulli, distribusi tersebut menjadi berbentuk lonceng.
Dalam Bab~\ref{chp 9}, kita akan membahas suatu teorema sangat umum yang disebut {\em
Teorema Limit Pusat} dan yang akan menjelaskan fenomena ini.
\end{example}

\begin{example}
Salah satu metode terkenal untuk menilai tangan bridge\index{bridge} adalah sebagai
berikut: sebuah as diberi nilai~4, raja~3, ratu~2, dan jack~1.  Semua kartu lainnya diberi
nilai~0.  \emx {Jumlah poin}\index{jumlah poin} tangan tersebut adalah jumlah nilai
kartu-kartu dalam tangan.  (Sebenarnya perhitungannya lebih rumit daripada ini karena
memperhitungkan ketiadaan kartu pada kembang tertentu dan sebagainya, tetapi di sini kita
meninjau bentuk jumlah poin yang disederhanakan ini.)  Jika sebuah kartu dibagikan secara
acak kepada seorang pemain, jumlah poin kartu tersebut memiliki distribusi
$$
p_X = \pmatrix{
0 & 1 & 2 & 3 & 4 \cr
36/52 & 4/52 & 4/52 & 4/52 & 4/52\cr}.
$$

Mari kita pandang seluruh tangan yang terdiri atas 13~kartu sebagai 13 percobaan saling
bebas dengan distribusi bersama ini.  (Sekali lagi, hal ini tidak sepenuhnya tepat karena
di sini kita menganggap bahwa kita selalu memilih kartu dari tumpukan lengkap.)  Maka
distribusi jumlah poin $C$ untuk tangan tersebut dapat diperoleh dari program {\bf
NFoldConvolution}\index{NFoldConvolution (program)} dengan menggunakan distribusi satu
kartu dan memilih $n = 13$.  Pemain dengan jumlah poin~13 atau lebih dikatakan memiliki
\emx {penawaran pembuka.} Peluang memiliki penawaran pembuka adalah
$$
P(C \geq 13)\ .
$$

Karena kita memiliki distribusi $C$, peluang ini mudah dihitung.  Dengan menghitungnya,
kita memperoleh
$$
P(C \geq 13) = .2845\ ,
$$
sehingga menurut model yang disederhanakan ini, kira-kira satu dari empat tangan seharusnya
memiliki penawaran pembuka.  Pembahasan yang lebih realistis mengenai masalah ini dapat
ditemukan dalam buku Epstein,\index{EPSTEIN, R.} \emx {The Theory of Gambling and Statistical Logic.}\footnote{R.~A.
Epstein,  \emx {The Theory of Gambling and Statistical Logic,} rev.~ed.\ (New
York: Academic Press, 1977).}
\end{example}
\par
Untuk distribusi khusus tertentu, kita dapat menemukan ungkapan bagi distribusi yang
dihasilkan dengan mengonvolusikan distribusi tersebut dengan dirinya sendiri sebanyak
$n$~kali.
\par
Konvolusi dua distribusi binomial,\index{konvolusi!distribusi binomial} yang satu dengan
parameter $m$ dan $p$ serta yang lain dengan parameter $n$ dan $p$, merupakan distribusi
binomial dengan parameter $(m+n)$ dan $p$.  Fakta ini mudah diperoleh dengan meninjau
percobaan yang terdiri atas pelemparan koin mula-mula sebanyak $m$ kali, kemudian sebanyak
$n$ kali lagi.
\par
Konvolusi $k$ distribusi geometrik\index{konvolusi!distribusi geometrik} dengan parameter
yang sama, yaitu $p$, merupakan distribusi binomial negatif dengan parameter $p$ dan $k$.  Hal ini
dapat dilihat dengan meninjau percobaan yang terdiri atas pelemparan koin sampai sisi
kepala ke-$k$ muncul.


\exercises 
\begin{LJSItem}


\i\label{exer 7.1.1} Sebuah dadu dilempar tiga kali.  Carilah peluang bahwa jumlah
hasilnya
\begin{enumerate}
\item lebih besar daripada 9.

\item merupakan bilangan ganjil.
\end{enumerate}

\i\label{exer 7.1.2} Harga suatu saham pada hari perdagangan tertentu berubah menurut
distribusi
$$
p_X = \pmatrix{
-1 & 0 & 1 & 2 \cr
1/4 & 1/2 & 1/8 & 1/8\cr}.$$
Carilah distribusi perubahan harga saham setelah dua hari perdagangan (dengan perubahan
yang saling bebas).

\i\label{exer 7.1.3} Misalkan $X_1$ dan $X_2$ peubah acak yang saling bebas dengan
distribusi yang sama,
$$
p_X = \pmatrix{
0 & 1 & 2 \cr
1/8 & 3/8 & 1/2\cr}.$$
Carilah distribusi jumlah $X_1 + X_2$.

\i\label{exer 7.1.4} Dalam satu permainan tertentu, Anda memenangkan uang sebesar $X$ dengan distribusi
$$
p_X = \pmatrix{
1 & 2 & 3 \cr
1/4 & 1/4 & 1/2\cr}.$$
Dengan menggunakan program {\bf NFoldConvolution}, carilah distribusi total kemenangan
Anda setelah sepuluh permainan (yang saling bebas).  Gambarkan distribusi ini.

\i\label{exer 7.1.5} Tinjaulah dua percobaan berikut: percobaan pertama memiliki hasil $X$
yang bernilai~0,~1, dan~2 dengan peluang sama; percobaan kedua menghasilkan hasil $Y$
(yang saling bebas dari hasil pertama) bernilai~3 dengan peluang 1/4 dan~4 dengan peluang
3/4.  Carilah distribusi
\begin{enumerate}
\item $Y + X$.

\item $Y - X$.
\end{enumerate}

\i\label{exer 7.1.6} Orang-orang tiba di suatu antrean menurut skema berikut: Dalam setiap
menit, sebanyak 0~atau~1 orang tiba.  Peluang 1 orang tiba adalah $p$, sedangkan peluang
tidak seorang pun tiba adalah $q = 1 - p$.  Misalkan $C_r$ adalah banyaknya pelanggan
yang tiba dalam $r$~menit pertama.  Tinjaulah proses percobaan Bernoulli dengan
keberhasilan jika seseorang tiba dalam satu satuan waktu dan kegagalan jika tidak seorang
pun tiba dalam satu satuan waktu.  Misalkan $T_r$ adalah banyaknya kegagalan sebelum
keberhasilan ke-$r$.
\begin{enumerate}
\item Tentukan distribusi $T_r$.

\item Tentukan distribusi $C_r$.

\item Carilah rata-rata dan varians banyaknya pelanggan yang tiba dalam $r$ menit
pertama.
\end{enumerate}

\i\label{exer 7.1.7} 
\begin{enumerate}
\item Sebuah dadu dilempar tiga kali dengan hasil $X_1$,~$X_2$, dan~$X_3$.  Misalkan
$Y_3$ adalah nilai maksimum yang diperoleh.  Tunjukkan bahwa
$$
P(Y_3 \leq j) = P(X_1 \leq j)^3\ .
$$
Gunakan hasil ini untuk mencari distribusi~$Y_3$.  Apakah $Y_3$ memiliki distribusi
berbentuk lonceng?
\item
Sekarang misalkan $Y_n$ adalah nilai maksimum ketika $n$ dadu dilempar.  Carilah
distribusi $Y_n$.  Apakah distribusi ini berbentuk lonceng untuk nilai $n$ yang besar?
\end{enumerate}

\i\label{exer 7.1.10} Seorang pemain bisbol akan bermain dalam World Series.  Berdasarkan
permainannya selama musim tersebut, Anda memperkirakan bahwa jika ia mendapat giliran
memukul empat kali dalam satu pertandingan, banyaknya pukulan berhasil yang diperolehnya
memiliki distribusi
$$
p_X = \pmatrix{
0 & 1 & 2 & 3 & 4 \cr
.4 & .2 & .2 & .1 & .1\cr}.$$

Anggaplah pemain tersebut mendapat giliran memukul empat kali dalam setiap pertandingan
seri.
\begin{enumerate}
\item
Misalkan $X$ menyatakan banyaknya pukulan berhasil yang ia peroleh dalam satu seri.
Dengan menggunakan program {\bf NFoldConvolution}, carilah distribusi $X$ untuk setiap
panjang seri yang mungkin: empat pertandingan, lima pertandingan, enam pertandingan,
tujuh pertandingan.

\item
Dengan menggunakan salah satu distribusi yang diperoleh dalam bagian (a), carilah peluang
rata-rata pukulannya melebihi .400 dalam seri empat pertandingan.  (Rata-rata pukulan
adalah banyaknya pukulan berhasil dibagi banyaknya giliran memukul.)

\item Diberikan distribusi $p_X$, berapakah rata-rata pukulan jangka panjangnya?
\end{enumerate}

\i\label{exer 7.1.11} Buktikan bahwa Anda tidak dapat memberi bobot pada dua dadu
sedemikian rupa sehingga peluang setiap jumlah dari 2~sampai~12 sama.  (Pastikan Anda
meninjau kasus ketika satu sisi atau lebih muncul dengan peluang nol.)

\i\label{exer 7.1.12} (L\'evy\footnote{Lihat M. Krasner dan B. Ranulae, ``Sur une Propriet\'e
des Polynomes de la Division du Circle"; dan catatan berikut oleh J.
Hadamard berikut, dalam \emx {C.\ R.\ Acad.\ Sci.,} vol.~204 (1937), pp.~397--399.})
Anggaplah $n$ suatu bilangan bulat yang bukan prima.  Tunjukkan bahwa Anda dapat mencari
dua distribusi $a$~dan~$b$ pada bilangan bulat tak-negatif sedemikian sehingga konvolusi
$a$~dan~$b$ merupakan distribusi seragam pada himpunan 0, 1, 2, \dots, $n - 1$.  Jika
$n$ prima, hal ini tidak mungkin, tetapi buktinya tidak begitu mudah.  (Anggaplah baik
$a$ maupun $b$ tidak terkonsentrasi pada 0.)

\i\label{exer 7.1.13} Anggaplah Anda bermain craps dengan dadu yang diberi bobot sebagai
berikut: sisi dua, tiga, empat, dan lima semuanya muncul dengan peluang yang sama,
$(1/6) + r$.  Sisi satu dan enam muncul dengan peluang $(1/6) - 2r$, dengan
$0 < r < .02$.  Tulislah program komputer untuk mencari peluang menang dalam craps
dengan dadu ini, lalu gunakan program Anda untuk mencari nilai~$r$ yang membuat craps
menjadi permainan yang menguntungkan pemain dengan dadu tersebut.
\end{LJSItem}




\choice{}{\section{Jumlah Peubah Acak Kontinu}\label{sec 7.2}
Dalam bagian ini kita meninjau versi kontinu dari masalah yang diajukan pada bagian
sebelumnya: Bagaimanakah distribusi jumlah peubah acak yang saling bebas?

\subsection*{Konvolusi}
\begin{definition} Misalkan $X$ dan $Y$ dua peubah acak kontinu dengan fungsi densitas
masing-masing $f(x)$ dan $g(y)$.  Anggaplah $f(x)$ dan $g(y)$ keduanya didefinisikan
untuk semua bilangan real.  Maka \emx {konvolusi}\index{konvolusi} $f*g$ dari
$f$~dan~$g$ adalah fungsi yang diberikan oleh
\begin{eqnarray*} (f*g)(z) &=& \int_{-\infty}^{+\infty} f(z - y) g(y)\,dy \\
         &=& \int_{-\infty}^{+\infty} g(z - x) f(x)\, dx\ .
\end{eqnarray*} 
\end{definition} 

Definisi ini sejalan dengan definisi konvolusi dua fungsi distribusi yang diberikan dalam
Bagian~\ref{sec 7.1}.  Jadi, tidak mengherankan bahwa jika $X$ dan $Y$ saling bebas,
densitas jumlah keduanya merupakan konvolusi densitas masing-masing.  Fakta ini dinyatakan
sebagai teorema di bawah, sedangkan buktinya diserahkan sebagai latihan (lihat
Latihan~\ref{exer 7.2.0.5}).

\begin{theorem} Misalkan $X$ dan $Y$ dua peubah acak yang saling bebas dengan fungsi
densitas $f_X(x)$ dan $f_Y(y)$ yang didefinisikan untuk semua~$x$.  Maka jumlah
$Z = X + Y$ merupakan peubah acak dengan fungsi densitas $f_Z(z)$, dengan $f_Z$
merupakan konvolusi $f_X$~dan~$f_Y$.
\end{theorem}
\par
Untuk memahami hasil penting ini dengan lebih baik, kita akan melihat beberapa contoh.
\pagebreak[4]
\subsection*{Jumlah Dua Peubah Acak Seragam yang Saling Bebas}
\begin{example}\label{exam 7.6}
Misalkan kita memilih secara acak dan saling bebas dua bilangan dari interval $[0,1]$
dengan densitas peluang seragam\index{konvolusi!densitas seragam}.  Tentukan densitas
jumlah keduanya.

Misalkan $X$ dan $Y$ peubah acak yang menjelaskan pilihan kita dan $Z = X + Y$ jumlah
keduanya.  Maka kita memiliki

$$
f_X(x) = f_Y(x) = \left \{ \begin{array}{ll}
                               1 & \;\mbox{jika $0 \leq x \leq 1$,} \\
                               0 & \;\mbox{selainnya;}
                  \end{array}
         \right. 
$$
dan fungsi densitas jumlahnya diberikan oleh
$$
f_Z(z) = \int_{-\infty}^{+\infty} f_X(z - y) f_Y(y)\,dy\ .
$$
Karena $f_Y(y) = 1$ jika $0 \leq y \leq 1$ dan 0 selainnya, bentuk ini menjadi
$$
f_Z(z) = \int_0^1 f_X(z - y)\,dy\ .
$$
Sekarang integrannya bernilai 0 kecuali jika $0 \leq z - y \leq 1$ (yaitu, kecuali jika $z - 1 \leq y
\leq z$), dan dalam hal itu nilainya~1.  Jadi, jika $0 \leq z \leq 1$,
kita memiliki
$$
f_Z(z) = \int_0^z \, dy = z\ ,
$$
sedangkan jika $1 < z \leq 2$, kita memiliki
$$
f_Z(z) = \int_{z - 1}^1\, dy = 2 - z\ ,
$$
dan jika $z < 0$ atau $z > 2$, kita memiliki $f_Z(z) = 0$ (lihat
Gambar~\ref{fig 7.5}).  Jadi,
$$
f_Z(z) = \left \{ \begin{array}{ll}
                               z,          & \;\mbox{jika $0 \leq z \leq 1,$} \\
                               2-z,        & \;\mbox{jika $1 < z \leq 2,$} \\
                               0,          & \;\mbox{selainnya.}
\end{array}
\right. 
$$
\putfig{3.5truein}{PSfig7-5}{Konvolusi dua densitas seragam.}{fig 7.5} %4.5truein
Perhatikan bahwa hasil ini sesuai dengan hasil Contoh~\ref{exam 2.1.4.5}.
\end{example}


\subsection*{Jumlah Dua Peubah Acak Eksponensial yang Saling Bebas}
\begin{example}\label{exam 7.7}
Misalkan kita memilih secara acak dua bilangan dari interval $[0,\infty)$ dengan densitas
\emx {eksponensial}\index{konvolusi!densitas eksponensial} berparameter $\lambda$.
Tentukan densitas jumlah keduanya.

Misalkan $X$, $Y$, dan $Z = X + Y$ menyatakan peubah-peubah acak yang bersangkutan, dan
$f_X$,~$f_Y$, serta~$f_Z$ densitasnya.  Maka
$$
f_X(x) = f_Y(x) = \left \{ \begin{array}{ll}
           \lambda e^{-\lambda x}, & \;\mbox{jika $x \geq 0$},\\
                                 0, & \;\mbox{selainnya;}
                  \end{array}
         \right. 
$$
dan dengan demikian, jika $z > 0$,
\begin{eqnarray*}
f_Z(z) &=& \int_{-\infty}^{+\infty} f_X(z - y) f_Y(y)\, dy \\
       &=& \int_0^z \lambda e^{-\lambda(z - y)} \lambda e^{-\lambda y}\, dy \\
       &=& \int_0^z \lambda^2 e^{-\lambda z}\, dy \\
       &=& \lambda^2 z e^{-\lambda z},\\
\end{eqnarray*}
sedangkan jika $z < 0$, $f_Z(z) = 0$ (lihat Gambar~\ref{fig 7.6}).  Jadi,
\putfig{3.5truein}{PSfig7-6}
{Konvolusi dua densitas eksponensial dengan $\lambda = 1$.}{fig 7.6} 
$$
f_Z(z) = \left \{ \begin{array}{ll}
                       \lambda^2 z e^{-\lambda z}, 
                                          & \;\mbox{jika $z \geq 0$},\\
                                      0,  & \;\mbox{selainnya.}
                  \end{array}
         \right. 
$$
\end{example}

\subsection*{Jumlah Dua Peubah Acak Normal yang Saling Bebas}
\begin{example}\label{exam 7.8}
Merupakan fakta yang menarik dan penting bahwa konvolusi dua densitas normal dengan
rata-rata $\mu_1$~dan~$\mu_2$ serta varians $\sigma_1$~dan~$\sigma_2$ kembali berupa
densitas normal, dengan rata-rata $\mu_1 + \mu_2$ dan varians $\sigma_1^2 + \sigma_2^2$.
Kita akan menunjukkan hal ini dalam kasus khusus ketika kedua peubah acak berdistribusi
normal standar.  Kasus umum dapat dikerjakan dengan cara yang sama, tetapi perhitungannya
lebih rumit.  Cara lain untuk menunjukkan hasil umum diberikan dalam
Contoh~\ref{exam 10.3.3}.
\par
Misalkan $X$ dan $Y$ dua peubah acak yang saling bebas, masing-masing dengan densitas
{\em normal} standar\index{konvolusi!densitas normal} (lihat
Contoh~\ref{exam 5.16}).  Kita memiliki
$$
f_X(x) = f_Y(y) = \frac 1{\sqrt{2\pi}} e^{-x^2/2}\ ,
$$
dan dengan demikian
\begin{eqnarray*}
f_Z(z) &=& f_X * f_Y(z) \\
&=& \frac 1{2\pi} \int_{-\infty}^{+\infty} e^{-(z -
y)^2/2} e^{-y^2/2}\, dy \\
       &=& \frac 1{2\pi} e^{-z^2/4} \int_{-\infty}^{+\infty} e^{-(y - z/2)^2}\,
dy \\
       &=& \frac 1{2\pi} e^{-z^2/4}\sqrt {\pi} \biggl[\frac 1{\sqrt {\pi}}\int_{-\infty}^\infty
e^{-(y-z/2)^2}\,dy\ \biggr]\ .\\
\end{eqnarray*}
Ungkapan di dalam kurung siku sama dengan 1 karena merupakan integral fungsi densitas
normal dengan $\mu = 0$ dan $\sigma = \sqrt 2$.  Jadi, kita memiliki
$$
f_Z(z) = \frac 1{\sqrt{4\pi}} e^{-z^2/4}\ .
$$
\end{example}

\subsection*{Jumlah Dua Peubah Acak Cauchy yang Saling Bebas}
\begin{example}\label{exam 7.9}
Pilihlah secara acak dua bilangan dari interval $(-\infty,+\infty)$ dengan densitas
Cauchy\index{konvolusi!densitas Cauchy} berparameter $a = 1$ (lihat
Contoh~\ref{exam 5.20}).  Maka
$$
f_X(x) = f_Y(x) = \frac 1{\pi(1 + x^2)}\ ,
$$
dan $Z = X + Y$ memiliki densitas
$$
f_Z(z) = \frac 1{\pi^2} \int_{-\infty}^{+\infty} \frac {1}{1 + (z - y)^2} \frac
{1}{1 + y^2} \, dy\ .
$$
Integral ini memerlukan usaha tertentu, dan di sini kita hanya memberikan hasilnya
(lihat Bagian~\ref{sec 10.3}, atau Dwass\footnote{M. Dwass, ``On the Convolution of Cauchy
Distributions,"  \emx {American Mathematical Monthly,} vol.~92, no.~1, (1985),
pp.~55--57; lihat juga R.~Nelson, surat kepada Editor, ibid., p.~679.}):
$$
f_Z(z) = \frac {2}{\pi(4 + z^2)}\ .
$$
\par
Sekarang, misalkan kita mencari fungsi densitas \emx {rata-rata}
$$
A = (1/2)(X + Y)
$$ 
dari $X$~dan~$Y$.  Maka $A = (1/2)Z$.  Latihan~\ref{sec 5.2}.\ref{exer
5.2.18} menunjukkan bahwa jika $U$ dan $V$ dua peubah acak kontinu dengan fungsi densitas
masing-masing $f_U(x)$ dan $f_V(x)$, dan jika $V = aU$, maka
$$
f_V(x) = \biggl(\frac 1a\biggr)f_U\biggl(\frac xa\biggr)\ .
$$
Dengan demikian, kita memiliki
$$
f_A(z) = 2f_Z(2z) = \frac 1{\pi(1 + z^2)}\ .
$$
Jadi, rata-rata dua peubah acak yang masing-masing memiliki densitas Cauchy kembali
memiliki densitas Cauchy; sifat luar biasa ini
merupakan kekhasan densitas Cauchy.  Salah satu akibatnya ialah bahwa jika galat dalam
suatu proses pengukuran memiliki densitas Cauchy dan Anda merata-ratakan sejumlah
pengukuran, rata-rata tersebut tidak dapat diharapkan lebih akurat daripada salah satu
pengukuran Anda sendiri!
\end{example}

\subsection*{Densitas Rayleigh}\index{fungsi densitas!Rayleigh}\index{densitas
Rayleigh}
\begin{example}\label{exam 7.10}
Misalkan $X$ dan $Y$ dua peubah acak normal standar yang saling bebas.  Sekarang misalkan
kita menempatkan titik~$P$ pada bidang-$xy$ dengan koordinat $(X,Y)$ dan bertanya:
tentukan densitas kuadrat jarak~$P$ dari titik asal.  (Kita telah menyimulasikan masalah ini
dalam Contoh~\ref{exam 5.19}.)  Di sini, dengan notasi sebelumnya, kita memiliki
$$
f_X(x) = f_Y(x) = \frac 1{\sqrt{2\pi}} e^{-x^2/2}\ .
$$
Selain itu, jika $X^2$ menyatakan kuadrat $X$, maka (lihat Teorema~\ref{thm 5.1} dan
pembahasan setelahnya)

\begin{eqnarray*}
f_{X^2}(r) &=& \left \{ \begin{array}{ll}
        \frac{1}{2\sqrt r} (f_X(\sqrt r) + f_X(-\sqrt r))  & \;\mbox{jika $r > 0,$} \\
                                     0                     & \;\mbox{selainnya.}
\end{array}
\right. \\
           &=& \left \{ \begin{array}{ll}
        \frac{1}{\sqrt {2 \pi r}} (e^{-r/2}) \hspace{.8in} & \;\mbox{jika $r > 0,$} \\
                                       0                   & \;\mbox{selainnya.}
\end{array}
\right. \\
\end{eqnarray*}
Ini merupakan densitas gamma dengan $\lambda = 1/2$, $\beta = 1/2$ (lihat
Contoh~\ref{exam 7.7}).  Sekarang misalkan $R^2 = X^2 + Y^2$.  Maka
\begin{eqnarray*}
f_{R^2}(r) &=& \int_{-\infty}^{+\infty} f_{X^2}(r - s) f_{Y^2}(s)\, ds \\
           &=& \frac 1{4\pi} \int_{-\infty}^{+\infty} e^{-(r - s)/2} 
{\frac{r-s}{2}}^{-1/2} e^{-s} {\frac{s}{2}}^{-1/2}\, ds\ , \\
           &=& \left \{\begin{array}{ll}
                       {\frac {1}{2}} e^{-r^2/2}, & \;\mbox{jika $r \geq 0,$} \\
                               0,                 & \;\mbox{selainnya.}
                  \end{array}
         \right. 
\end{eqnarray*}
Jadi, $R^2$ memiliki densitas gamma dengan $\lambda = 1/2$, $\beta = 1$.  Hasil ini dapat
kita tafsirkan sebagai densitas kuadrat jarak~$P$ dari pusat suatu sasaran jika
koordinatnya berdistribusi normal.

Densitas peubah acak $R$ diperoleh dari densitas $R^2$ dengan cara biasa (lihat
Teorema~\ref{thm 5.1}), dan kita memperoleh
$$
f_R(r) = \left \{ \begin{array}{ll}
                       \frac 12 e^{-r^2/2} \cdot 2r = re^{-r^2/2}, 
                                          & \;\mbox{jika $r \geq 0,$} \\
                               0,         & \;\mbox{selainnya.}
                  \end{array}
         \right. 
$$

Para fisikawan akan mengenali fungsi ini sebagai densitas Rayleigh.  Hasil kita di sini
sesuai dengan simulasi dalam Contoh~\ref{exam 5.19}.
\end{example}

\subsection*{Densitas Khi-Kuadrat}\index{densitas khi-kuadrat}\index{fungsi densitas!khi-kuadrat}
Secara lebih umum, metode yang sama menunjukkan bahwa jumlah kuadrat~$n$ peubah acak
yang saling bebas dan berdistribusi normal dengan rata-rata~0 serta simpangan baku~1
memiliki densitas gamma dengan $\lambda = 1/2$ dan $\beta = n/2$.  Densitas semacam itu
disebut \emx {densitas khi-kuadrat} dengan $n$ derajat kebebasan.  Densitas ini
diperkenalkan dalam Bab~\ref{chp 5}.  Dalam Contoh~\ref{exam 5.20}, kita menggunakan
densitas ini untuk menguji hipotesis bahwa dua sifat saling bebas.
\par
Penggunaan penting lain dari densitas khi-kuadrat adalah membandingkan data eksperimen
dengan distribusi diskret teoretis untuk melihat apakah data mendukung model teoretis
tersebut.  Secara lebih khusus, misalkan kita memiliki percobaan dengan himpunan hasil
berhingga.  Jika himpunan hasilnya terhitung, kita mengelompokkannya menjadi sejumlah
berhingga himpunan hasil.  Kita mengusulkan distribusi teoretis yang menurut kita akan
memodelkan percobaan dengan baik.  Kita memperoleh data dengan mengulang percobaan
beberapa kali.  Sekarang kita ingin memeriksa seberapa baik distribusi teoretis tersebut
cocok dengan data.
\par
Misalkan $X$ peubah acak yang mewakili hasil teoretis dalam model percobaan, dan misalkan
$m(x)$ fungsi distribusi $X$.  Dengan cara yang serupa dengan yang dilakukan dalam
Contoh~\ref{exam 5.20}, kita menghitung nilai ungkapan
$$
V = \sum_x \frac{(o_x - n \cdot m(x))^2}{n \cdot m(x)}\ ,
$$
dengan penjumlahan dilakukan atas semua hasil yang mungkin $x$, $n$ adalah banyaknya
titik data, dan $o_x$ menyatakan banyaknya hasil bertipe $x$ yang diamati dalam data.
Untuk nilai $n$ sedang atau besar, kuantitas $V$ kemudian berdistribusi khi-kuadrat
secara hampiran dengan $\nu - 1$ derajat kebebasan, dengan $\nu$ menyatakan banyaknya
hasil yang mungkin.  Bukti pernyataan ini berada di luar cakupan buku, tetapi kita akan
menggambarkan kewajaran pernyataan tersebut dalam contoh berikutnya.  Jika nilai $V$
sangat besar dibandingkan dengan fungsi densitas khi-kuadrat yang sesuai, kita cenderung
menolak hipotesis bahwa model tersebut sesuai bagi percobaan yang sedang dihadapi.
Sekarang kita memberikan contoh prosedur ini.

\begin{example}
Misalkan kita diberi sebuah dadu.  Kita ingin menguji hipotesis bahwa dadu tersebut
seimbang.  Jadi, distribusi teoretis kita adalah distribusi seragam pada bilangan bulat
di antara 1 dan 6.  Jika kita melempar dadu $n$ kali, banyaknya titik data yang diharapkan
untuk setiap tipe adalah $n/6$.  Jadi, jika $o_i$ menyatakan banyaknya titik data aktual
bertipe $i$, untuk $1 \le i \le 6$, maka ungkapan
$$V = \sum_{i = 1}^6 \frac{(o_i - n/6)^2}{n/6}$$
berdistribusi khi-kuadrat secara hampiran dengan 5 derajat kebebasan.
\par
Sekarang misalkan kita benar-benar melempar dadu 60 kali dan memperoleh data dalam
Tabel~\ref{table 7.1}.
\begin{table}
\centering
\begin{tabular}{|c|c|}
\hline
 Hasil & Frekuensi Teramati \\
\hline 1 & 15\\
\hline 2 & \hspace{.08in}8 \\
\hline 3 & \hspace{.08in}7 \\
\hline 4 & \hspace{.08in}5 \\
\hline 5 & \hspace{.08in}7 \\
\hline 6 & 18 \\
\hline
\end{tabular}
\caption{Data teramati.}
\label{table 7.1}
\end{table}
Jika kita menghitung $V$ untuk data ini, kita memperoleh nilai 13.6.  Grafik densitas
khi-kuadrat dengan 5 derajat kebebasan ditunjukkan pada Gambar~\ref{fig 7.7}.  Terlihat
bahwa nilai sebesar 13.6 jarang diperoleh oleh $V$ jika dadu seimbang, sehingga kita akan
menolak hipotesis bahwa dadu tersebut seimbang.  (Ketika menggunakan uji ini, ahli
statistika akan menolak hipotesis jika data memberikan nilai $V$ yang lebih besar daripada
95\% nilai yang diharapkan diperoleh jika hipotesis tersebut benar.)
\putfig{3.5truein}{PSfig7-7}{Densitas khi-kuadrat dengan 5 derajat kebebasan.}{fig 7.7} 
\par
Pada Gambar~\ref{fig 7.8}, kita menampilkan hasil pelemparan dadu 60 kali, kemudian
menghitung $V$, lalu mengulang percobaan ini 1000 kali.  Program yang melakukan
perhitungan ini disebut {\bf DieTest}.\index{DieTest (program)}  Kita menumpangkan
densitas khi-kuadrat dengan 5 derajat kebebasan; terlihat bahwa nilai data cukup sesuai
dengan kurva, yang mendukung pernyataan bahwa densitas khi-kuadrat merupakan densitas yang
tepat untuk digunakan.
\putfig{4.5truein}{PSfig7-8}{Pelemparan dadu seimbang.}{fig 7.8} 
\end{example}

Sejauh ini kita telah melihat beberapa kasus khusus penting ketika integral konvolusi
dapat dihitung secara eksplisit.  Secara umum, konvolusi dua densitas kontinu tidak dapat
dihitung secara eksplisit, sehingga kita harus menggunakan metode numerik.  Untungnya,
metode tersebut ternyata sangat efektif, setidaknya bagi densitas terbatas.

\subsection*{Percobaan Saling Bebas}
Sekarang kita meninjau secara singkat distribusi jumlah~$n$ peubah acak yang saling bebas,
yang semuanya memiliki fungsi densitas yang sama.  Jika $X_1$,~$X_2$, \dots,~$X_n$
adalah peubah-peubah acak tersebut dan $S_n = X_1 + X_2 +\cdots+ X_n$ adalah jumlahnya,
maka kita akan memiliki
$$
f_{S_n}(x) = \left( f_{X_1} * f_{X_2} *\cdots* f_{X_n} \right)(x)\ ,
$$
dengan ruas kanan merupakan konvolusi lipat-$n$.  Dalam kasus sederhana tertentu,
densitas ini dapat dihitung untuk nilai umum~$n$.

\begin{example}\label{exam 7.12}
Misalkan $X_i$ berdistribusi seragam\index{konvolusi!densitas seragam} pada interval
$[0,1]$.  Maka
$$
f_{X_i}(x) = \left \{ \begin{array}{ll}
                         1,         & \;\mbox{jika $0 \leq x \leq 1,$} \\
                         0,         & \;\mbox{selainnya,}
                   \end{array}
          \right. 
$$
dan $f_{S_n}(x)$ diberikan oleh rumus\index{USPENSKY, J. B.}\footnote{J.~B. Uspensky,
{\em Introduction to Mathematical Probability} (New York: McGraw-Hill, 1937),
p.~277.}
$$
f_{S_n}(x) = \left \{ \begin{array}{ll}
                         \frac 1{(n - 1)!} \sum_{0 \leq j \leq x} (-1)^j
 { n \choose j} (x - j)^{n - 1},     & \;\mbox{jika $0 < x < n,$} \\       
                         0,          & \;\mbox{selainnya.}
                   \end{array}
          \right. 
$$
Densitas $f_{S_n}(x)$ untuk~$n = 2$,~4, 6, 8,~10 ditunjukkan pada
Gambar~\ref{fig 7.9}.

\putfig{4.5truein}{PSfig7-9}{Konvolusi $n$ densitas seragam.}{fig 7.9} 


Jika $X_i$ berdistribusi normal,\index{konvolusi!densitas normal standar} dengan
rata-rata~0 dan varians~1, maka (bdk.~Contoh~\ref{exam 7.8})
$$
f_{X_i}(x) = \frac 1{\sqrt{2\pi}} e^{-x^2/2}\ ,
$$
dan
$$
f_{S_n}(x) = \frac 1{\sqrt{2\pi n}} e^{-x^2/2n}\ .
$$
Di sini densitas $f_{S_n}$ untuk~$n = 5$,~10, 15, 20,~25 ditunjukkan pada
Gambar~\ref{fig 7.10}.

\putfig{4.5truein}{PSfig7-10}{Konvolusi $n$ densitas normal standar.}{fig 7.10} 

Jika semua $X_i$ berdistribusi eksponensial,\index{konvolusi!densitas eksponensial}
dengan rata-rata~$1/\lambda$, maka
$$
f_{X_i}(x) = \lambda e^{-\lambda x}\ ,
$$
dan
$$
f_{S_n}(x) = \frac {\lambda e^{-\lambda x}(\lambda x)^{n - 1}}{(n - 1)!}\ .
$$
Dalam kasus ini densitas $f_{S_n}$ untuk~$n = 2$,~4, 6, 8,~10 ditunjukkan pada
Gambar~\ref{fig 7.11}.
\putfig{4.5truein}{PSfig7-11}
{Konvolusi $n$ densitas eksponensial dengan $\lambda = 1$.}{fig 7.11} 
\end{example}


\exercises
\begin{LJSItem}

\i\label{exer 7.2.0.5}  Misalkan $X$ dan $Y$ peubah acak bernilai real yang saling bebas 
dengan fungsi densitas masing-masing $f_X(x)$ dan $f_Y(y)$.  Tunjukkan bahwa fungsi densitas 
jumlah $X + Y$ merupakan konvolusi fungsi $f_X(x)$ dan $f_Y(y)$.
\emx {Petunjuk}:  Misalkan $\bar X$ peubah acak gabungan $(X, Y)$.  Maka fungsi densitas
gabungan $\bar X$ ialah $f_X(x)f_Y(y)$ karena $X$ dan $Y$ saling bebas.  Sekarang
hitung peluang bahwa $X+Y \le z$ dengan mengintegrasikan fungsi densitas gabungan pada
daerah yang sesuai di bidang.  Ini menghasilkan fungsi distribusi kumulatif $Z$. 
Lalu diferensiasikan fungsi ini terhadap $z$ untuk memperoleh fungsi densitas sebagai fungsi $z$.

\i\label{exer 7.2.1} Misalkan $X$ dan $Y$ peubah acak yang saling bebas dan didefinisikan
pada ruang $\Omega$, dengan fungsi densitas masing-masing $f_X$~dan~$f_Y$. 
Misalkan $Z = X + Y$.  Tentukan densitas $f_Z$ untuk~$Z$ jika
\begin{enumerate}
\item $$f_X(x) = f_Y(x) = \left \{ \begin{array}{ll}
                              1/2,   & \;\mbox{jika $-1 \leq x \leq +1,$} \\
                              0,     & \;\mbox{selainnya.}
                  \end{array}
         \right. $$

 
\item $$f_X(x) = f_Y(x) = \left \{ \begin{array}{ll}
                              1/2,   & \;\mbox{jika $3 \leq x \leq 5,$} \\
                              0,     & \;\mbox{selainnya.}
                  \end{array}
         \right. $$


\item $$f_X(x) = \left \{ \begin{array}{ll}
                              1/2,   & \;\mbox{jika $-1 \leq x \leq 1,$} \\
                              0,     & \;\mbox{selainnya.}
                  \end{array}
         \right. $$
\smallskip
$$f_Y(x) = \left \{ \begin{array}{ll}
                              1/2,   & \;\mbox{jika $3 \leq x \leq 5,$} \\
                              0,     & \;\mbox{selainnya.}
                  \end{array}
         \right. $$
\item Apa yang dapat Anda simpulkan tentang himpunan $E = \{\,z : f_Z(z)> 0\,\}$ pada setiap
kasus?
\end{enumerate}

\i\label{exer 7.2.2} Misalkan lagi bahwa $Z = X + Y$.  Tentukan $f_Z$ jika
\begin{enumerate}
\item $$f_X(x) = f_Y(x) = \left \{ \begin{array}{ll} 
                    x/2, & \mbox{jika $0 < x < 2,$} \\
                      0, & \mbox{selainnya}.
                  \end{array}
         \right. $$

\item $$f_X(x) = f_Y(x) = \left \{ \begin{array}{ll} 
                 (1/2)(x - 3), & \mbox{jika $3 < x < 5,$} \\
                            0, & \mbox{selainnya}. 
                  \end{array}
         \right. $$

\item $$f_X(x) = \left \{ \begin{array}{ll} 
                       1/2,    & \mbox{jika $0 < x < 2,$} \\
                            0, & \mbox{selainnya},
                  \end{array}
         \right. $$
\smallskip
 $$f_Y(x) = \left \{ \begin{array}{ll} 
                          x/2, & \mbox{jika $0 < x < 2,$} \\
                            0, & \mbox{selainnya}.
                  \end{array}
         \right. $$
\item Apa yang dapat Anda simpulkan tentang himpunan $E = \{\,z : f_Z(z)> 0\,\}$ pada setiap
kasus?
\end{enumerate}

\i\label{exer 7.2.3} Misalkan $X$, $Y$, dan $Z$ peubah acak yang saling bebas
dengan 
$$f_X(x) = f_Y(x) = f_Z(x)  = \left \{ \begin{array}{ll}
                              1,     & \mbox{jika $0 < x < 1,$} \\
                              0,     & \mbox{selainnya.}
                  \end{array}
         \right. $$
\noindent Misalkan $W = X + Y + Z$.  Tentukan $f_W$ secara langsung dan bandingkan jawaban Anda
dengan hasil dari rumus dalam Contoh~\ref{exam 7.12}.   \emx {Petunjuk}:  Lihat 
Contoh~\ref{exam 7.6}.

\i\label{exer 7.2.3.5} Misalkan $X$ dan $Y$ saling bebas dan $Z = X + Y$.  Tentukan
$f_Z$ jika

\begin{enumerate}
\item $$f_X(x) = \left \{ \begin{array}{ll}
                              \lambda e^{-\lambda x}, & \mbox{jika $x > 0,$} \\
                               0,                     & \mbox{selainnya.}
                  \end{array}
         \right. $$
\smallskip
$$f_Y(x) = \left \{ \begin{array}{ll}
                              \mu e^{-\mu x}, & \mbox{jika $x > 0,$} \\
                              0,              & \mbox{selainnya.}
                  \end{array}
         \right. $$
\item $$\ \ \ f_X(x) = \left \{ \begin{array}{ll}
                              \lambda e^{-\lambda x}, & \mbox{jika $x > 0,$} \\
                              0,                      & \mbox{selainnya.}
                  \end{array}
         \right. $$

\smallskip
$$f_Y(x) = \left \{ \begin{array}{ll}
                    1, & \mbox{jika $0 < x < 1,$} \\
                    0, & \mbox{selainnya.}
                  \end{array}
         \right. $$
\end{enumerate}

\i\label{exer 7.2.100} Misalkan lagi bahwa $Z = X + Y$.  Tentukan $f_Z$ jika
\begin{eqnarray*}
f_X(x) &=& \frac 1{\sqrt{2\pi}\sigma_1} e^{-(x - \mu_1)^2/2\sigma_1^2} \\
f_Y(x) &=& \frac 1{\sqrt{2\pi}\sigma_2} e^{-(x - \mu_2)^2/2\sigma_2^2}\ .
\end{eqnarray*}

%********Soal berikut terlalu sulit.
\istar\label{exer 7.2.101} Misalkan $R^2 = X^2 + Y^2$.  Tentukan $f_{R^2}$ dan $f_R$ jika
\begin{eqnarray*}
f_X(x) &=& \frac 1{\sqrt{2\pi}\sigma_1} e^{-(x - \mu_1)^2/2\sigma_1^2} \\
f_Y(x) &=& \frac 1{\sqrt{2\pi}\sigma_2} e^{-(x - \mu_2)^2/2\sigma_2^2}\ .
\end{eqnarray*}

\i\label{exer 7.2.101.5} Misalkan $R^2 = X^2 + Y^2$.  Tentukan $f_{R^2}$ dan $f_R$ jika
$$
f_X(x) = f_Y(x) = \left \{ \begin{array}{ll}
                    1/2, & \mbox{jika $-1 \leq x \leq 1,$} \\
                    0,   & \mbox{selainnya.}
                  \end{array}
         \right. 
$$

\i\label{exer 7.2.102} Asumsikan bahwa waktu pelayanan seorang nasabah di bank
berdistribusi eksponensial dengan waktu pelayanan rata-rata 2~menit.  Misalkan $X$ adalah total waktu
pelayanan bagi 10 nasabah.  Perkirakan peluang bahwa $X > 22$ menit.

\i\label{exer 7.2.9} Misalkan $X_1$,~$X_2$, \dots,~$X_n$ adalah $n$ peubah acak yang saling
bebas, masing-masing dengan densitas eksponensial yang memiliki rata-rata~$\mu$.  Misalkan $M$ adalah
nilai \emx {minimum} dari $X_j$.  Tunjukkan bahwa densitas~$M$ bersifat
eksponensial dengan rata-rata $\mu/n$.   \emx {Petunjuk}:  Gunakan fungsi distribusi kumulatif.

\i\label{exer 7.2.103} Sebuah perusahaan membeli 100 bohlam yang masing-masing memiliki umur
pakai berdistribusi eksponensial dengan rata-rata 1000 jam.  Berapa lama waktu yang diharapkan hingga bohlam pertama
putus?  (Lihat Latihan~\ref{exer 7.2.9}.)

\i\label{exer 7.2.104} Sebuah perusahaan asuransi mengasumsikan bahwa selang waktu antarklaim pada
setiap polis pemilik rumahnya berdistribusi eksponensial dengan rata-rata~$\mu$.  Perusahaan itu
ingin memperkirakan $\mu$ dengan merata-ratakan selang waktu untuk sejumlah polis,
tetapi cara ini kurang praktis karena selang waktu antarklaim sekitar
30~tahun.  Atas saran Galambos\index{GALAMBOS, J.}\footnote{J. Galambos,  \emx {Introductory
Probability Theory} (New York: Marcel Dekker, 1984), p.~159.}, perusahaan membagi
nasabahnya ke dalam kelompok beranggotakan~50 dan mengamati waktu terjadinya klaim pertama dalam
setiap kelompok.  Tunjukkan bahwa cara ini praktis untuk memperkirakan nilai
~$\mu$.

\i\label{exer 7.2.105} Partikel mengalami tumbukan yang menyebabkannya terbelah menjadi
dua bagian, dengan setiap bagian merupakan suatu fraksi dari partikel induk.  Misalkan fraksi ini
berdistribusi seragam antara 0~dan~1.  Dengan mengikuti satu partikel melalui
beberapa pembelahan, kita memperoleh fraksi partikel semula $Z_n = X_1
\cdot X_2 \cdot\dots\cdot X_n$, dengan setiap $X_j$ berdistribusi seragam
antara 0~dan~1.  Tunjukkan bahwa densitas peubah acak $Z_n$ ialah
$$
f_n(z) = \frac 1{(n - 1)!}( -\log z)^{n - 1}.
$$
 \emx {Petunjuk}: Tunjukkan bahwa $Y_k = -\log X_k$ berdistribusi eksponensial.  Gunakan
hasil ini untuk mencari fungsi densitas $S_n = Y_1 + Y_2 +\cdots+ Y_n$, lalu
gunakan hasil tersebut untuk memperoleh distribusi kumulatif dan densitas $Z_n = e^{-S_n}$.

\i\label{exer 7.2.106} Asumsikan bahwa $X_1$ dan $X_2$ peubah acak yang saling bebas, masing-masing
memiliki densitas eksponensial berparameter~$\lambda$.  Tunjukkan bahwa $Z = X_1 - X_2$ memiliki
densitas
$$
f_Z(z) = (1/2)\lambda e^{-\lambda |z|}\ .
$$

\i\label{exer 7.2.107} Misalkan kita ingin menguji apakah sebuah koin seimbang.  Kita melempar koin $n$
kali dan mencatat banyaknya kemunculan sisi ekor sebagai $X_0$ serta
banyaknya kemunculan sisi kepala sebagai $X_1 = n - X_0$.  Sekarang kita tetapkan
$$
Z= \sum_{i = 0}^1 \frac {(X_i - n/2)^2}{n/2}\ .
$$
Maka, untuk koin seimbang, $Z$ kira-kira berdistribusi khi-kuadrat dengan $2 -
1 = 1$ derajat kebebasan.  Verifikasikan ini melalui simulasi komputer, mula-mula untuk koin
seimbang ($p~=~1/2$), lalu untuk koin tak seimbang ($p~=~1/3$).

\i\label{exer 7.2.108} Verifikasikan jawaban Anda untuk Latihan~\ref{exer 7.2.1}(a) melalui
simulasi komputer: pilih $X$ dan $Y$ dari $[-1,1]$ dengan densitas seragam, lalu hitung
$Z = X + Y$.  Ulangi percobaan ini 500 kali dan catat hasilnya dalam diagram
batang pada $[-2,2]$ dengan 40~batang.  Apakah densitas $f_Z$ yang dihitung dalam Latihan~\ref{exer
7.2.1}(a) menggambarkan bentuk diagram batang Anda?  Cobalah juga untuk Latihan~
\ref{exer 7.2.1}(b)~dan~Latihan~\ref{exer 7.2.1}(c).

\i\label{exer 7.2.109} Verifikasikan jawaban Anda untuk Latihan~\ref{exer 7.2.2} melalui
simulasi komputer.

\i\label{exer 7.2.110} Verifikasikan jawaban Anda untuk Latihan~\ref{exer 7.2.3} melalui
simulasi komputer.

\i\label{exer 7.2.18} \emx {Dukungan} fungsi $f(x)$ didefinisikan sebagai himpunan
$$
\{x\ :\ f(x) > 0\}\ .$$
Misalkan $X$ dan $Y$ dua peubah acak kontinu dengan
fungsi densitas masing-masing $f_X(x)$ dan $f_Y(y)$, dan misalkan dukungan kedua
fungsi densitas itu masing-masing interval $[a, b]$ dan $[c, d]$.  Tentukan dukungan
fungsi densitas peubah acak $X+Y$.

\i\label{exer 7.2.111} Misalkan $X_1$,~$X_2$, \dots,~$X_n$ barisan peubah acak yang saling
bebas dan semuanya memiliki fungsi densitas yang sama, $f_X$, dengan dukungan $[a,b]$ (lihat
Latihan~\ref{exer 7.2.18}).  Misalkan $S_n = X_1 + X_2 +\cdots+ X_n$, dengan fungsi
densitas $f_{S_n}$.  Tunjukkan bahwa dukungan~$f_{S_n}$ adalah interval
$[na,nb]$.   \emx {Petunjuk}: Tulis $f_{S_n} = f_{S_{n - 1}} * f_X$.  Kemudian gunakan
Latihan~\ref{exer 7.2.18} untuk membuktikan hasil yang diminta dengan induksi.

\i\label{exer 7.2.112} Misalkan $X_1$,~$X_2$, \dots,~$X_n$ barisan peubah acak yang saling
bebas dan semuanya memiliki fungsi densitas yang sama, $f_X$.  Misalkan $A = S_n/n$ merupakan
rata-rata peubah-peubah tersebut.  Tentukan $f_A$ jika
\begin{enumerate}
\item $f_X(x) = (1/\sqrt{2\pi}) e^{-x^2/2}$ (densitas normal).

\item $f_X(x) = e^{-x}$ (densitas eksponensial).
\par
\noindent \emx {Petunjuk}: Nyatakan $f_A(x)$ dalam bentuk $f_{S_n}(x)$.
\end{enumerate}

\end{LJSItem}}

